\documentclass[11pt,a4paper]{article}
\usepackage[utf8]{inputenc}
\usepackage[T1]{fontenc}
\usepackage{lmodern}
\usepackage{xcolor}
\renewcommand{\contentsname}{Spis treści}
\renewcommand{\tablename}{Tabela}
\hyphenpenalty=3000 \tolerance=2000 \emergencystretch=2em
\usepackage[margin=2.1cm]{geometry}
\usepackage{booktabs}
\usepackage{longtable}
\usepackage{enumitem}
\usepackage{titlesec}
\usepackage{parskip}
\titleformat{\section}{\large\bfseries}{\thesection.}{0.6em}{}
\titleformat{\subsection}{\normalsize\bfseries}{\thesubsection.}{0.6em}{}
\setlist{nosep,leftmargin=1.4em}
\newcommand{\ohm}{$\Omega$}

\title{\textbf{AMPERE POINT --- custom device}\\[4pt]
\large PROPOZYCJA v2: pomiary wspierające diagnostykę przypadków\\
,,u klienta nie działa, u nas działa'' (moduł A)\\[2pt]
\normalsize wersja 2 --- do zatwierdzenia; zastępuje propozycję v1 po uwagach z 2026-07-09}
\author{Opracowanie wewnętrzne AMPERE POINT}
\date{9 lipca 2026}

\begin{document}
\maketitle
\vspace{-1.4em}

\section{Uwagi do v1 --- przyjęte}
P1 (konformancja CP) i P3 (funkcje L--N/PE) --- producent testuje to u siebie, przełączanie stanów
w praktyce serwisowej nie zawodzi; P2 --- Tuya to poglądowa informacja dla klienta, nie przyrząd,
rozliczenia nie opierają się na liczniku ładowarki, a ewentualna korekta wymagałaby wielopoziomowej
pracy po stronie producenta. \textbf{Wniosek:} nie testujemy zgodności produktu --- wspieramy
\emph{proces diagnostyki}, zwłaszcza najdroższy jego przypadek: \textbf{reklamacja działa u klienta,
sprzęt na stole działa poprawnie}.

\section{Skąd się bierze ,,u klienta nie działa, u nas działa''}
Ładowarka jest sprawna --- różni się \emph{środowisko}. Typowe różnice (na bazie materiałów
w folderze: instrukcja serii P, katalog usterek, przewodnik pomiarów, doświadczenia serwisu):
\begin{enumerate}
\item \textbf{Zasilanie klienta}: napięcie za wysokie (sieć z PV w południe, $>$253~V) lub za
niskie; \textbf{miękka sieć} --- napięcie siada pod prądem ładowania (długi WLZ, przedłużacz,
słabe gniazdo); zapady od rozruchów innych odbiorników; asymetria; krótkie przerwy.
\item \textbf{Instalacja klienta}: RCD wyzwala, bo \emph{suma} upływów (inne odbiory $+$ auto $+$
ładowarka) przekracza 30~mA, choć każdy z osobna jest OK; zły typ RCD; grzejące się gniazdo;
zamiana L--N / słabe PE (ładowarka poprawnie blokuje --- klient widzi ,,nie działa'').
\item \textbf{Auto klienta}: nietypowe zachowania CP po stronie pojazdu (usypianie/wybudzanie,
harmonogramy auta, wolne przejścia stanów, chwilowe upływy DC z OBC wyzwalające RDC-DD 6~mA),
przejściówki i kable.
\item \textbf{Łączność}: słabe WiFi w miejscu montażu $\to$ ,,appka nie działa'' $=$ reklamacja,
choć elektrycznie wszystko gra.
\item \textbf{Środowisko}: temperatura (wstrzymania termiczne), wilgoć/kondensacja (upływy nocą).
\end{enumerate}
Diagnoza wymaga: \textbf{(a) zarejestrować incydent tam, gdzie występuje} (u klienta, w chwili
awarii), \textbf{(b) odtworzyć te warunki na stole}, aż sprawny egzemplarz zawiedzie tak samo.
Stąd trzy tryby pracy modułu A + analityka.

\section{Proponowane pakiety D1--D5}
\subsection{D1 --- tryb REJESTRATOR: czarna skrzynka u klienta (dni--tygodnie)}
Moduł A wpięty \textbf{przelotowo} w zasilanie ładowarki klienta (przenośne P/B/Q: gniazdo
klienta $\to$ moduł A $\to$ wtyczka ładowarki; wallbox: przez rozłączkę serwisową). Rejestruje
z rozdzielczością półokresu: U RMS (min/max, zapady/przepięcia ze znacznikiem czasu), prąd, moc,
\textbf{upływ różnicowy sumaryczny toru} (RCM --- kluczowe przy ,,RCD u klienta wyzwala''),
impedancję źródła $\Delta$U/$\Delta$I liczoną z naturalnych zmian obciążenia, temperatury,
\textbf{RSSI WiFi i zajętość kanału} w miejscu montażu (radio ESP32 --- za darmo), zdarzenia
(zanik L, wstrzymania). Log na SD $+$ zdalny podgląd (HA), praca bez nadzoru. Wynik: \emph{dowód},
co dokładnie działo się w chwili awarii.
\subsection{D2 --- tryb PRZELOT auto$\leftrightarrow$ładowarka: podsłuch prawdziwej sesji}
Nowy kabel pojazdowy z modułu A (wtyk Type~2 ,,do auta''): ładowarka klienta wpięta w -X1, auto
w kabel --- moduł A siedzi \textbf{w środku sesji} i biernie rejestruje: przebieg CP obu stron
(stany, duty, anomalie czasowe), PP, U/I/P (przez istniejące -B1/-T/-U2), upływ różnicowy pary
auto$+$ładowarka. Rozstrzyga spory ,,wina auta czy ładowarki'' i łapie przypadki zależne od
konkretnego pojazdu. (Podsłuch wysokoimpedancyjny --- bez ingerencji w sesję.)
\subsection{D3 --- tryb SYMULATOR: odtwarzanie warunków klienta na stole}
W sekcji bench-test (zasilanie DUT z panelu): \textbf{wtrącana impedancja linii} (odczepy
0,15/0,3/0,6~\ohm, rezystory drutowe z termikami --- symulacja miękkiej sieci/przedłużacza),
\textbf{regulowany upływ tła} L$\to$PE przed DUT (rozszerzenie istniejącego generatora -A3
o przełącznik strony --- symulacja ,,doładowanych'' upływów instalacji), \textbf{krótkie przerwy
i zapady} (szybki stycznik z wyzwalaniem czasowym), zamiana L--N / przerwa PE (przekaźniki ---
tu w roli \emph{reprodukcji zgłoszenia}, nie testu zgodności). Opcja etapu 2: autotransformator
16~A (używany) dla pełnej regulacji 190--260~V. Procedura ,,\textbf{replay}'': firmware odtwarza
sekwencję zdarzeń zarejestrowaną w D1 aż do powtórzenia usterki.
\subsection{D4 --- diagnostyka łączności (zero sprzętu)}
Pomiar na miejscu u klienta (moduł A lub sam ESP32 serwisowy): mapa RSSI 2,4~GHz, zajętość kanałów,
stabilność połączenia z chmurą Tuya (opóźnienia/utraty), test ,,ładowarka widzi router, ale nie
chmurę''. Duża część zgłoszeń ,,nie działa'' to appka --- ten pakiet zamyka je w minuty.
\subsection{D5 --- analityka i pamięć przypadków}
Raport diagnostyczny z D1--D4 w jednym formacie; \textbf{biblioteka profili klienckich} (nagrane
warunki: ,,miękka sieć 0,6~\ohm'', ,,upływ tła 18~mA'', ,,WiFi $-$82~dBm'') do regresji nowych
partii importu; powiązanie z katalogiem usterek U-xxx (nowe wpisy z dowodami z rejestratora).

\section{Wpływ na konstrukcję modułu A (krótko; do zatwierdzenia)}
\begin{longtable}{@{}p{0.07\textwidth}p{0.55\textwidth}p{0.14\textwidth}p{0.14\textwidth}@{}}
\toprule
\textbf{Pak.} & \textbf{Zmiana konstrukcyjna} & \textbf{Koszt} & \textbf{Praca}\\
\midrule
D1 & tor przelotowy 1F/3F-16 A w sekcji bench-test: wejście CEE16$+$Schuko $\to$ \textbf{druga
wiązka przez okna istniejących -B1 i -T1} (tryb R $=$ tor główny bez prądu, więc czujniki wolne)
$\to$ gniazdo wyjściowe; przekaźnik -K34 przełącza dzielniki -U2 na tor przelotowy; zasilanie
logiki jak dotąd z osobnego przewodu & 150--300 zł & 1 dzień\\
D2 & kabel z wtykiem pojazdowym Type~2 $+$ złącze panelowe; odczep przelotowy CP/PP z -X1
z przekaźnikiem trybu (podsłuch zamiast PM701E); blokada programowa: w trybie P stycznik -K1
otwarty & 250--450 zł & 0,5--1 dzień\\
D3 & w sekcji bench-test: 3 odczepy rezystancji linii (drutowe $+$ termiki do pętli -F11),
2 przekaźniki (L--N, PE), szybki stycznik przerw, przełącznik strony dla -A3; \emph{opcja:}
autotransformator 16 A używany & 200--400 zł (opcja $+$400--700) & 1 dzień\\
D4 & bez zmian sprzętowych (radio ESP32 $+$ firmware) & 0 & 0,5 dnia\\
D5 & firmware $+$ szablon raportu $+$ format profili na SD & 0 & 1 dzień\\
\midrule
 & \textbf{Razem (bez opcji)} & \textbf{600--1150 zł} & \textbf{4--4,5 dnia}\\
\bottomrule
\end{longtable}
Bez zmian: tor mocy E1 (tryby R/P używają go biernie), łańcuch E2, szafa B, sekwencja pomiarów
obowiązkowych. Warunek: sekcja bench-test (decyzja Z3) --- D1/D3 jej wymagają. Dokumentacyjnie:
nowy arkusz E4 (tor przelotowy, symulator, przełączniki trybów) --- po zatwierdzeniu.

\section{Przebieg diagnostyki po zmianie (docelowy proces)}
zgłoszenie $\to$ \textbf{D4} zdalnie/na miejscu (łączność? --- 30\% przypadków kończy się tu)
$\to$ \textbf{D1} rejestrator u klienta na 2--7 dni (incydent złapany z dowodem) $\to$ analiza
$\to$ \textbf{D3} replay na stole na sprawnym egzemplarzu $\to$ wniosek: wada środowiska (raport
dla klienta/elektryka) albo wada ukryta (raport z dowodami do producenta) $\to$ wpis do katalogu
U-xxx $+$ profil do biblioteki (\textbf{D5}) $\to$ regresja przyszłych partii tym profilem.

\section{Decyzja}
\textbf{Z20 (v2):} zakres D1--D5, budżet $\sim$0,6--1,15 tys. zł ($+$opcja autotransformatora),
$\sim$4--4,5 dnia pracy; D1/D3 warunkowane zatwierdzeniem sekcji bench-test (Z3). Rekomendowana
kolejność wdrożenia: \textbf{D4 $\to$ D1 $\to$ D5 $\to$ D3 $\to$ D2}.

\end{document}
